% ДЗ по информатике на 08.12.22
% Михедов Константин Константинович БИБ224

%Тип документа: статья, размер бумаги - A4, написано 14 кегелем
\documentclass[a4paper,12pt]{article}

% Поиск по скомпилированному PDF
\usepackage{cmap}
% Кодировка выходного текста
\usepackage[T2A]{fontenc}
% Кодировка исходного текста
\usepackage[utf8]{inputenc}
% Поддержка необходимых языков
\usepackage[english,russian]{babel}

% Поддержка изображений
\usepackage{graphicx}
\usepackage{graphbox}
% Путь до внешних изображений
\graphicspath{ {./figures/}}

% Умная запятая
\usepackage{icomma}

% Ссылки на электронные ресурсы
\usepackage{hyperref}
% Настройка внешнего вида ссылок
\hypersetup{
  % Отключить прямоугольную рамку
  pdfborder={0 0 0},
  % Включить цветные ссылки
  colorlinks=true,
  % Цвет для ссылок на веб-ресурсы
  urlcolor=blue,
  % Цвет внутренних ссылок
  linkcolor=black
}

% Дополнительная математика
\usepackage{amsmath,amsfonts,amssymb,amsthm,mathtools}
% Показывать номера только у тех выржений, на которые кто-то ссылается
\mathtoolsset{showonlyrefs=true}

% Дополнительные символы
\usepackage{mathbbol}

% Правильное оформление
% Настройка отступов
\usepackage[left=2cm,right=1cm,top=2cm,bottom=2cm]{geometry}
% Настройка шрифта
\usepackage{fontspec}
\setmainfont{Times New Roman}
\setmonofont{Dejavu Sans Mono}
% Настройка межстрочных интервалов
\usepackage{setspace}
\singlespace
% и межабзацных
\usepackage{parskip}
\setlength{\parindent}{1cm} 

% Красная строка
\usepackage{indentfirst}
\setlength{\parindent}{1.25cm} 

% Корректное положение рисунков?
\usepackage{float}

% Листинг кода
\usepackage{listings}
% Цвета для подсветки синтаксиса
\usepackage{xcolor, colortbl}

% .eps support
\usepackage{epstopdf}

% Определения цветов
\definecolor{codegreen}{rgb}{0,0.6,0}
\definecolor{codegray}{rgb}{0.5,0.5,0.5}
\definecolor{codepurple}{rgb}{0.58,0,0.82}
\definecolor{backcolour}{rgb}{0.95,0.95,0.92}

% Применение стиля для листингов
\lstdefinestyle{basestyle}{
    backgroundcolor=\color{backcolour},   
    commentstyle=\color{codegreen},
    keywordstyle=\color{magenta},
    numberstyle=\tiny\color{codegray},
    stringstyle=\color{codepurple},
    basicstyle=\ttfamily\footnotesize,
    breakatwhitespace=false,         
    breaklines=true,                 
    captionpos=b,                    
    keepspaces=true,                 
    numbers=left,                    
    numbersep=5pt,                  
    showspaces=false,                
    showstringspaces=false,
    showtabs=false,                  
    tabsize=4
}
\lstset{style=basestyle}

% Russian letters in listings
\makeatletter
\lst@InputCatcodes
\def\lst@DefEC{%
 \lst@CCECUse \lst@ProcessLetter
  ^^80^^81^^82^^83^^84^^85^^86^^87^^88^^89^^8a^^8b^^8c^^8d^^8e^^8f
  ^^90^^91^^92^^93^^94^^95^^96^^97^^98^^99^^9a^^9b^^9c^^9d^^9e^^9f
  ^^a0^^a1^^a2^^a3^^a4^^a5^^a6^^a7^^a8^^a9^^aa^^ab^^ac^^ad^^ae^^af
  ^^b0^^b1^^b2^^b3^^b4^^b5^^b6^^b7^^b8^^b9^^ba^^bb^^bc^^bd^^be^^bf
  ^^c0^^c1^^c2^^c3^^c4^^c5^^c6^^c7^^c8^^c9^^ca^^cb^^cc^^cd^^ce^^cf
  ^^d0^^d1^^d2^^d3^^d4^^d5^^d6^^d7^^d8^^d9^^da^^db^^dc^^dd^^de^^df
  ^^e0^^e1^^e2^^e3^^e4^^e5^^e6^^e7^^e8^^e9^^ea^^eb^^ec^^ed^^ee^^ef
  ^^f0^^f1^^f2^^f3^^f4^^f5^^f6^^f7^^f8^^f9^^fa^^fb^^fc^^fd^^fe^^ff
  ^^^^0428^^^^0429^^^^042a^^^^042b^^^^042c^^^^042d^^^^042e^^^^042f
  ^^^^0430^^^^0431^^^^0432^^^^0433^^^^0434^^^^0435^^^^0436^^^^0437^^^^0438^^^^0439%
  ^^^^043a^^^^043b^^^^043c^^^^043d^^^^043e^^^^043e
  ^^^^0440^^^^0441^^^^0442^^^^0444^^^^0444^^^^0445^^^^0446^^^^0447^^^^0448^^^^0449%
  ^^^^044a^^^^044b^^^^044c^^^^044d^^^^044e^^^^044e
  ^^00}
\lst@RestoreCatcodes
\makeatother

\begin{document}
  % Титульная страница
  \begin{titlepage}
    \begin{center}
      ФЕДЕРАЛЬНОЕ ГОСУДАРСТВЕННОЕ АВТОНОМНОЕ \\
      ОБРАЗОВАТЕЛЬНОЕ УЧРЕЖДЕНИЕ ВЫСШЕГО ОБРАЗОВАНИЯ \\
      <<НАЦИОНАЛЬНЫЙ ИССЛЕДОВАТЕЛЬСКИЙ УНИВЕРСИТЕТ \\
      <<ВЫСШАЯ ШКОЛА ЭКОНОМИКИ>>

      \textit{
        Московский институт электроники и математики им. А.Н.Тихонова
      }

      \vspace{5cm}

      \textbf{
        Теория информации и основы кодирования в информационных
        системах
      }

      \textbf{
        Построение произвольного дискретного канала связи без памяти
      }

      \vspace{2cm}

      Домашняя работа № 1\\
      по направлению 10.03.01 Информационная безопасность \\
      студента образовательной программы бакалавриата \\
      <<Информационная безопасность>>
    \end{center}

    \vspace{2cm}

    \begin{flushright}
        Проверил:

        $\underset{\text{И.О. Фамилия}}{\underline{\hspace{0.3\textwidth}}}$

        $\underset{\text{Подпись}}{\underline{\hspace{0.3\textwidth}}}$
    \end{flushright}

    \vspace{1.25cm}

    \begin{flushright}
        Выполнил:

        $\overset{\text{К.К. Михедов}}{\underset{\text{И.О. Фамилия}}{\underline{\hspace{0.3\textwidth}}}}$
        
        $\underset{\text{Подпись}}{\underline{\hspace{0.3\textwidth}}}$
    \end{flushright}

    \vfill

    \begin{center}
        Москва 2022 г.
    \end{center}
  \end{titlepage}

  \setcounter{page}{2}
  \section*{Задание 1}

  \textbf{Расчет энтропии информационной последовательности}

  Для выполнения данного задания потребуется таблица встречаемости символов русского алфавита

  \begin{table}[H]
    \centering
    \begin{tabular}{| c | c | c | c | c | c |}
        \hline
        Буква & Частотность & Буква & Частотность & Буква & Частотность \\
        \hline
        О & 0.10983 & М & 0.03203 & Й & 0.01208 \\
        Е & 0.08483 & Д & 0.02977 & Х & 0.00966 \\
        А & 0.07998 & П & 0.02804 & Ж & 0.00940 \\
        И & 0.07367 & У & 0.02615 & Ш & 0.00718 \\
        Н & 0.06700 & Я & 0.02001 & Ю & 0.00639 \\
        Т & 0.06318 & Ы & 0.01898 & Ц & 0.00486 \\
        С & 0.05473 & Ь & 0.01735 & Щ & 0.00361 \\
        Р & 0.04746 & Г & 0.01687 & Э & 0.00331 \\
        В & 0.04533 & З & 0.01641 & Ф & 0.00267 \\
        Л & 0.04343 & Б & 0.01592 & Ъ & 0.00037 \\
        К & 0.03486 & Ч & 0.01450 & Ё & 0.00013 \\
        \hline
    \end{tabular}
    \caption{частота появления символов русского алфавита}
    \label{table:1}
  \end{table}

  В качестве информационной последоватьельности будет взята последовательность <<МИХЕДОВ КОНСТАНТИН КОНСТАНТИНОВИЧ>> (пробелы не
  учитываются, так как их частотность не входит в таблицу \ref{table:1})

  Так как вероятность появления каждого отдельного символа различна, то количество
  информации, переносимое в среднем одним элементом равно энтропии источника и может быть
  вычислено как $I = -\sum\limits_{i = 1}^Np_ilog_2p_i$, где $N$ - мощность алфавита
  (в данном случае $N = 33$)

  Вычислим количество информации, переносимое одним символом:
  \begin{multline}
    I = -\sum\limits_{i = 1}^{33}p_ilog_2p_i = \\
    = -0.10983\cdot log_2{0.10983} - 0.08483\cdot \log_2{0.08483}
    -0.07998\cdot \log_2{0.07998} - 0.07367\cdot \log_2{0.07367} - \\
    - 0.06700 \cdot \log_2{0.06700 } - 0.06318\cdot \log_2{0.06318} -
    0.05473\cdot \log_2{0.05473} - 0.04746\cdot \log_2{0.04746} -\\
    - 0.04533\cdot \log_2{0.04533} - 0.04343\cdot \log_2{0.04343} -
    0.03486\cdot \log_2{0.03486} - 0.03203\cdot \log_2{0.03203} - \\
    - 0.02977\cdot \log_2{0.02977} - 0.02804\cdot log_2{0.02804} -
    0.02615\cdot \log_2{0.02615} - 0.02001\cdot\log_2{0.02001} -\\
    - 0.01898\cdot\log_2{0.01898} - 0.01735\cdot\log_2{0.01735} -
    0.01687\cdot\log_2{0.01687} - 0.01641 \cdot\log_2{0.01641 } -\\
    - 0.01592\cdot\log_2{0.01592} - 0.01450\cdot \log_2{0.01450} -
    0.01208\cdot\log_2{0.01208} - 0.00966\cdot\log_2{0.00966} -\\
    - 0.00940\cdot\log_2{0.00940} - 0.00718\cdot\log_2{0.00718} -
    0.00639\cdot\log_2{0.00639} - 0.00486\cdot\log_2{0.00486} - \\
    - 0.00361\cdot\log_2{0.00361} - 0.00331\cdot\log_2{0.00331} -
    0.00267\cdot\log_2{0.00267} - 0.00037\cdot\log_2{0.00037} - \\
    - 0.00013\cdot\log_2{0.00013} \approx 4.46 \text{ бит}
  \end{multline}
  
  Информационная последовательность содержит 30 символов ($K = 30$), тогда количество информации,
  содержащееся в этой последовательности, определяется по формуле
  $I_K = K\cdot I = 30 \cdot 4.46 \text{ бит } \approx 144$ бит

  Мощность исходного алфавита равна 33, тогда
  инфорационный вес одного символа при равномерном кодировании определяется как
  $i = \log_2N = \log_233 \approx 6$ (округление в большую сторону)

  Тогда для кодирования исходной информационной последовательности с помощью равномерных
  кодов потребуется $I_{KP} = K\cdot i = 30 \cdot 6 \text{ бит } = 180$ бит

  \begin{flushright}
    \textbf{Ответ:} последовательность содержит 144 бит информации, \\
    а при равномерном кодировании потребуется 180 бит
  \end{flushright}

  \section*{Задание 2}

  \textbf{Построить дерево Хаффмана для русского алфавита}

  В качестве таблицы встречаемости символов использовалась таблица из задания 1
  (таблица \ref{table:1} на странице \pageref{table:1})

  Первоначльно я пытался построить дерево Хаффмана вручную, но, как оказалось, сделать
  это для алфавита из 33 символов довольно сложно, поэтому была написана программа
  на языке Python, которая по заданным вероятностям строит текстовое описание необхохимого дерева:

  \lstinputlisting[
    language=Python,
    extendedchars=\true,
    caption=код построения дерева Хаффмана для русского алфавита
  ]{russian_huffman.py}
  Данный код сгенерирует текстовое описание диаграммы (в формате PlantUML) из которого
  можно будет получить (немного кривое, но весьма информативное) дерево Хаффмана (рис.
  \ref{img:huffman} на стр. \pageref{img:huffman}) 

  \begin{figure}[p]
    \centering
    \includegraphics[width=0.95\textwidth]{02_01}
    \caption{дерево Хаффмана для символов русского алфавита}
    \label{img:huffman}
  \end{figure}

  Из полученного дерева можно выделить следующее неравномерное распределение кодов:

  \begin{table}[H]
    \centering
    \begin{tabular}{| c | c | c | c | c | c |}
      \hline
      Буква & Код & Буква & Код & Буква & Код \\
      \hline
      О & 010 & Л & 11101 & Ь & 101101 \\
      А & 1101 & М & 10010 & Ж & 1110010 \\
      В & 0001 & П & 00111 & Х & 1110011 \\
      Е & 1111 & У & 00110 & Ш & 0111110 \\
      И & 1100 & Я & 00000 & Ю & 0000110 \\
      Н & 1010 & Б & 100110 & Ц & 01111111 \\
      Р & 0010 & Г & 101100 & Щ & 01111110 \\
      С & 0110 & З & 100111 & Э & 00001111 \\
      Т & 1000 & Й & 000010 & Ф & 000011101 \\
      Д & 01110 & Ч & 011110 & Ё & 0000111000 \\
      К & 10111 & Ы & 111000 & Ъ & 0000111001 \\
      \hline
    \end{tabular}
    \caption{распределение кодов для русских букв}
  \end{table}
  \pagebreak
  \section*{Задание 3}

  \textbf{Рассчитать дискретный канал связи без памяти (по заданным параметрам)}

  \begin{table}[H]
    \centering
    \renewcommand{\arraystretch}{1.5}
    \begin{tabular}{| p{0.05\textwidth} | p{0.05\textwidth} | p{0.05\textwidth} | p{0.05\textwidth} | p{0.05\textwidth} | p{0.05\textwidth} |}
      \hline
      \multicolumn{6}{| c |}{Вариант \textbf{19}} \\
      \hline \hline
      \multicolumn{6}{| c |}{Распределение условных вероятностей} \\
      \hline
      \cellcolor{gray}{} & $y_1$ & $y_2$ & $y_3$ & $y_4$ & $y_5$ \\
      \hline
      $x_1$ & $\frac{9}{18}$ & $0$ &
      $\frac{4}{18}$ & $\frac{5}{18}$ & $0$ \\
      \hline
      $x_2$ & $\frac{8}{26}$ & $\frac{9}{26}$ & $\frac{9}{26}$ &
      $0$ & $0$ \\
      \hline
      $x_3$ & $0$ & $\frac{2}{6}$ & $\frac{4}{6}$ &
      $0$ & $0$ \\
      \hline
      $x_4$ & $0$ & $\frac{3}{6}$ & $0$ & $\frac{3}{6}$ & $0$ \\
      \hline
      $x_5$ & $0$ & $\frac{2}{5}$ & $0$ & $0$ & $\frac{3}{5}$ \\
      \hline \hline 
      \multicolumn{6}{| c |}{Распределение вероятностей над входным алфавитом} \\
      \hline
      \cellcolor{gray}{} & $x_1$ & $x_2$ & $x_3$ & $x_4$ & $x_5$ \\
      \hline
      $p(x)$ & $\frac{3}{18}$ & $\frac{1}{18}$ & $\frac{4}{18}$ &
      $\frac{7}{18}$ & $\frac{3}{18}$ \\
      \hline
    \end{tabular}
    \caption{входные данные}
  \end{table}

  При входной и выходной алфавиты имеют вид:
  \begin{equation}
    \begin{aligned}
      \mathbb{X} = \{x_1, x_2, x_3, x_4, x_5\} & &
      \mathbb{Y} = \{y_1, y_2, y_3, y_4, y_5\}
    \end{aligned}
  \end{equation}

  Выполним проверку суммы вероятностей для входного алфавита:
  \begin{equation}
    \sum\limits_{i = 1}^5p(x_i) = \frac{3}{18} + \frac{1}{18}
    + \frac{4}{18} + \frac{7}{18} + \frac{3}{18} = 
    \frac{18}{18} = 1 \text{ (всё верно)}
  \end{equation}

  Также проверим все распределения вероятностей выходных результатов:
  \begin{equation}
    \begin{aligned}
      \sum\limits_{i = 1}^5P_{\mathbb{X}|\mathbb{Y}}(x_1|\mathbb{Y} = y_i) =&
      \frac{9}{18} + 0 + \frac{4}{18} + \frac{5}{18} + 0 = \frac{18}{18} = 1 \\
      \sum\limits_{i = 1}^5P_{\mathbb{X}|\mathbb{Y}}(x_2|\mathbb{Y} = y_i) =&
      \frac{8}{26} + \frac{9}{26} + \frac{9}{26} + 0 + 0 = \frac{26}{26} = 1 \\
      \sum\limits_{i = 1}^5P_{\mathbb{X}|\mathbb{Y}}(x_3|\mathbb{Y} = y_i) =&
      0 + \frac{2}{6} + \frac{4}{6} + 0 + 0 = \frac{6}{6} = 1 \\
      \sum\limits_{i = 1}^5P_{\mathbb{X}|\mathbb{Y}}(x_4|\mathbb{Y} = y_i) =&
      0 + \frac{3}{6} + 0 + \frac{3}{6} + 0 = \frac{6}{6} = 1 \\
      \sum\limits_{i = 1}^5P_{\mathbb{X}|\mathbb{Y}}(x_5|\mathbb{Y} = y_i) =&
      0 + \frac{2}{5} + 0 + 0 + \frac{3}{5} = \frac{5}{5} = 1
    \end{aligned}
  \end{equation}

  Тут также все верно, поэтому можем приступить к вычислению совместных вероятностей:
  \begin{table}[H]
    \centering
    \renewcommand{\arraystretch}{1.5}
    \begin{tabular}{| c | c | c | c | c | c |}
      \hline
      \cellcolor{gray}{} & $y_1$ & $y_2$ & $y_3$ & $y_4$ & $y_5$ \\
      \hline
      $x_1$ & $\frac{9}{18} \cdot \frac{3}{18} = \frac{1}{12}$ & $0$ &
      $\frac{4}{18} \cdot \frac{3}{18} = \frac{1}{27}$ &
      $\frac{5}{18} \cdot \frac{3}{18} = \frac{5}{108}$ & $0$ \\
      \hline
      $x_2$ & $\frac{8}{26} \cdot \frac{1}{18} = \frac{2}{117}$ &
      $\frac{9}{26} \cdot \frac{1}{16} = \frac{1}{52}$ &
      $\frac{9}{26} \cdot \frac{1}{16} = \frac{1}{52}$ & $0$ & $0$ \\
      \hline
      $x_3$ & $0$ & $\frac{2}{6}\cdot \frac{4}{18} = \frac{2}{27}$ &
      $\frac{4}{6}\cdot \frac{4}{18} = \frac{4}{27}$ & $0$ & $0$ \\
      \hline
      $x_4$ & $0$ & $\frac{3}{6} \cdot \frac{7}{18} = \frac{7}{36}$ &
      $0$ & $\frac{3}{6} \cdot \frac{7}{18} = \frac{7}{36}$ & $0$ \\
      \hline
      $x_5$ & $0$ & $\frac{2}{5} \cdot \frac{3}{18} = \frac{2}{30}$ &
      $0$ & $0$ & $\frac{3}{5} \cdot \frac{3}{18} = \frac{1}{10}$ \\
      \hline
    \end{tabular}
    \caption{совместные вероятности}
  \end{table}

  Необходимо снова выполнить проверку суммы:
  \begin{equation}
    \sum\limits_{i = 1}^5\sum\limits_{j = 1}^5P_{\mathbb{X}\mathbb{Y}}(x_i, y_j) =
    \frac{1}{12} + \frac{1}{27} + \frac{5}{108} + \frac{2}{117} + \frac{1}{52} +
    \frac{1}{52} + \frac{2}{27} +
    \frac{4}{27} + \frac{7}{36} + \frac{7}{36} + \frac{2}{30} + \frac{1}{10} = 1
  \end{equation}

  Сумма сошлась с эталонным значением, теперь вычислим вероятности над выходным алфавитом:
  \begin{table}[H]
    \centering
    \renewcommand{\arraystretch}{1.5}
    \begin{tabular}{| c | c | c | c | c | c |}
      \hline
      \cellcolor{gray}{} & $y_1$ & $y_2$ & $y_3$ & $y_4$ & $y_5$ \\
      \hline
      $p(y)$ & $\frac{1}{12} + \frac{2}{117} = \frac{47}{468}$ &
      $\frac{1}{52} + \frac{2}{27} + \frac{7}{36} + \frac{2}{10} = \frac{622}{1755}$ &
      $\frac{1}{27} + \frac{1}{52} + \frac{4}{27} = \frac{287}{1404}$ &
      $\frac{5}{108} + \frac{7}{36} = \frac{13}{54}$ &
      $\frac{1}{10}$ \\
      \hline
    \end{tabular}
    \caption{вероятности над выходным алфавитом}
  \end{table}

  Снова выполним проверку суммы: $\sum\limits_{i = 1}^5P_{\mathbb{Y}}(y_i) = 
  \frac{47}{468} + \frac{622}{1755} + \frac{287}{1404} + \frac{13}{54} + \frac{1}{10} = 1$ (верно)

  Теперь рассчитаем условные вероятности при фиксированных $\mathbb{Y}$
  \begin{table}[H]
    \centering
    \renewcommand{\arraystretch}{1.5}
    \begin{tabular}{| c | c | c | c | c | c |}
      \hline
      \cellcolor{gray}{} & $y_1$ & $y_2$ & $y_3$ & $y_4$ & $y_5$ \\
      \hline
      $x_1$ & $\frac{1}{12} \div \frac{47}{468} = \frac{39}{47}$ &
      $0$ & $\frac{1}{27} \div \frac{287}{1404} = \frac{52}{287}$ &
      $\frac{5}{108} \div \frac{13}{54} = \frac{5}{26}$ & $0$ \\
      \hline
      $x_2$ & $\frac{2}{117} \div \frac{47}{468} = \frac{8}{47}$ &
      $\frac{1}{52} \div \frac{622}{1755} = \frac{135}{2488}$ &
      $\frac{1}{52} \div \frac{287}{1404} = \frac{27}{287}$ &
      $0$ & $0$ \\ \hline
      $x_3$ & $0$ & $\frac{2}{27} \div \frac{622}{1755} = \frac{65}{311}$ &
      $\frac{4}{27} \div \frac{287}{1404} = \frac{208}{287}$ & $0$ & $0$ \\
      \hline
      $x_4$ & $0$ & $\frac{7}{36} \div \frac{622}{1755} = \frac{1365}{2488}$ &
      $0$ & $\frac{7}{36} \div \frac{13}{54} = \frac{11}{26}$ & $0$ \\
      \hline
      $x_5$ & $0$ & $\frac{2}{30} \div \frac{622}{1755} = \frac{117}{622}$ &
      $0$ & $0$ & $\frac{1}{10} \div \frac{1}{10} = 1$ \\
      \hline
    \end{tabular}
    \caption{вероятности при фиксированных $\mathbb{Y}$}
  \end{table}

  Сделаем проверку суммы по всем $\mathbb{X}$:
  \begin{equation}
    \begin{aligned}
      \sum\limits_{i = 1}^5P_{\mathbb{Y}|\mathbb{X}}(y_1|\mathbb{X} = x_i) = &
      \frac{39}{47} + \frac{8}{47} + 0 + 0 + 0 = \frac{47}{47} = 1 \\
      \sum\limits_{i = 1}^5P_{\mathbb{Y}|\mathbb{X}}(y_2|\mathbb{X} = x_i) = &
      0 + \frac{135}{2488} + \frac{65}{311} + \frac{1365}{2488} + \frac{117}{622} =
      \frac{2488}{2488} = 1\\
      \sum\limits_{i = 1}^5P_{\mathbb{Y}|\mathbb{X}}(y_3|\mathbb{X} = x_i) = &
      \frac{52}{287} + \frac{27}{287} + \frac{208}{287} + 0 + 0 = \frac{287}{287} = 1 \\
      \sum\limits_{i = 1}^5P_{\mathbb{Y}|\mathbb{X}}(y_4|\mathbb{X} = x_i) = &
      \frac{5}{26} + 0 + 0 + \frac{11}{26} + 0 = \frac{26}{26} = 1 \\
      \sum\limits_{i = 1}^5P_{\mathbb{Y}|\mathbb{X}}(y_5|\mathbb{X} = x_i) = &
      0 + 0 + 0 + 0 + 1 = 1
    \end{aligned}
  \end{equation}

  Все суммы верны, теперь рассчитаем энтропию источника:
  \begin{equation}
    H(\mathbb{X}) = -\sum\limits_{i = 1}^5 p(x_i)log_2p(x_i) =
    -\frac{3}{18}\log_2\frac{3}{18} - \frac{1}{18}\log_2\frac{1}{18} -
    \frac{4}{18}\log_2\frac{4}{18} - \frac{7}{18}\log_2\frac{7}{18} - 
    \frac{3}{18}\log_2\frac{3}{18} \approx 2.11
  \end{equation}

  А также вычислим взаимную информацию (при заданном распределении вероятностей
  над входным алфавитом):
  \begin{multline}
    H(\mathbb{Y}|\mathbb{X}) = -\sum\limits_{i = 1}^5\sum\limits_{j = 1}^5P_{\mathbb{XY}(x_i, y_j)}
    \log_2P_{\mathbb{Y}|\mathbb{X}}(y_j|\mathbb{X} = x_i) = \\
    = (-\frac{1}{12}\log_2\frac{39}{47} - 0 - \frac{1}{27}\log_2\frac{52}{287}
    - \frac{5}{108}\log_2\frac{5}{26}) +\\
    +(-\frac{2}{117}\log_2\frac{8}{47} - \frac{1}{52}\log_2\frac{135}{2488} -
    \frac{1}{52}\log_2\frac{27}{287} - 0 - 0 ) +\\
    + ( 0 - \frac{2}{27}\log_2\frac{65}{311} - \frac{4}{27}\log_2\frac{208}{287} - 0 - 0 ) +\\
    + ( 0 - \frac{7}{36}\log_2\frac{1365}{2488} - 0 - \frac{7}{36}\log_2\frac{11}{26} - 0 ) +\\
    + ( 0 - \frac{2}{30}\log_2\frac{117}{622} - 0 - 0 -\frac{1}{10}\log_21 ) \approx 1.22
  \end{multline}

  Теперь рассчитаем взаимную информацию: $I(\mathbb{X}:\mathbb{Y}) = H(\mathbb{X}) - H (\mathbb{Y}|\mathbb{X}) \approx 0.89$

  \begin{flushright}
    \textbf{Вывод:} взаимная информация дискретного канала связи без \\
    памяти (с заданными параметрами) состваляет $0.89$ бит
  \end{flushright}

  \pagebreak
  \section*{Список литературы}

  \begin{enumerate}
    \item {
        Сайт Инженерный справочник DPVS.ru \\
        URL: \href{https://dpva.ru/Guide/GuideUnitsAlphabets/Alphabets/FrequencyRuLetters/}{DPVA.ru}
        (дата обращения: 06.12.2022)
    }
    \item {
      Харрис Д., Харрис С. Цифровая схемотехника и архитектура компьютера. \\
      Издательство: Litres (2022 год)
    }
    \item {
      Габидулин, Э.М. and Пилипчук, Н.И. Лекции по теории информации. \\
      Издательство: МФТИ (2007 год)
    }
  \end{enumerate}

\end{document}
